\subsection{Milestone 2 -- Identificazione di modelli accurati}

\subsubsection{Run preliminare: 27 combinazioni (SMOTE e wrapper disabilitati)}

I risultati riportati in questa sezione si riferiscono a un run preliminare con la
configurazione default: 3 classificatori $\times$ 3 strategie di feature selection $\times$
3 strategie di balancing = 27 combinazioni, con SMOTE e metodi wrapper disabilitati.

La feature selection (NONE, InfoGain, Spearman) produce risultati identici alla terza
cifra decimale per ogni coppia (Classificatore, Balancing): il fenomeno è discusso in
dettaglio nella sezione successiva. Le righe della tabella riportano i valori con
\textit{No FS} come rappresentativi di tutte e tre le strategie di feature selection.

\begin{table}[h!]
\centering
\small
\begin{tabular}{@{}llccccc@{}}
\toprule
\textbf{Classificatore} & \textbf{Balancing} & \textbf{Precision} & \textbf{Recall} &
\textbf{AUC} & \textbf{Kappa} & \textbf{NPofB20} \\
\midrule
RandomForest  & None          & 0,73 & 0,52 & 0,94 & 0,58 & 0,88 \\
RandomForest  & Undersampling & 0,31 & 0,86 & 0,91 & 0,37 & 0,80 \\
RandomForest  & Oversampling  & 0,61 & 0,66 & 0,94 & 0,59 & 0,87 \\
NaiveBayes    & None          & 0,35 & 0,29 & 0,79 & 0,26 & 0,54 \\
NaiveBayes    & Undersampling & 0,34 & 0,32 & 0,76 & 0,26 & 0,53 \\
NaiveBayes    & Oversampling  & 0,34 & 0,32 & 0,78 & 0,26 & 0,54 \\
IBk           & None          & 0,53 & 0,51 & 0,75 & 0,47 & 0,56 \\
IBk           & Undersampling & 0,26 & 0,81 & 0,79 & 0,30 & 0,57 \\
IBk           & Oversampling  & 0,47 & 0,58 & 0,77 & 0,46 & 0,61 \\
\bottomrule
\end{tabular}
\caption{Risultati del run preliminare (No FS; metriche arrotondate a 2 cifre decimali).
I valori con InfoGain e Spearman sono identici a quelli riportati per ogni riga.}
\end{table}

\subsubsection{Feature selection come no-op su questo dataset}

Per ogni coppia (Classificatore, Balancing), i valori con No FS, InfoGain e Spearman sono
identici alla terza cifra decimale. La causa è che \texttt{InfoGainAttributeEval} +
\texttt{Ranker(threshold=0.0)} mantiene tutti gli attributi con InfoGain $> 0$: le 19
feature di OpenJPA sono tutte informative rispetto alla classe \texttt{isBuggy}, nessuna
ha InfoGain esattamente zero. Analogamente, Spearman con soglia 0 mantiene tutto. Il
modello addestrato è quindi identico al caso No FS.

Questo non è un difetto della pipeline ma un risultato valido: \textbf{tutte le feature
portano informazione} su questo dataset. I metodi wrapper (ForwardSearch, BackwardSearch)
potrebbero trovare sottoinsiemi più compatti, ma i tempi di esecuzione (1--3 ore per
singola combinazione) ne hanno finora impedito l'utilizzo. I run futuri con i wrapper
permetteranno di verificare se esiste un sottoinsieme ottimale che supera il plateau
attuale.

\subsubsection{Kappa inferiore a 0,6: limite strutturale, non un difetto}

Il valore $0{,}6$ è la soglia standard Landis \& Koch (1977) per l'``accordo sostanziale''
in diagnostica medica. Tutti i valori di questo run sono sotto $0{,}6$ (il massimo è
RF + Oversampling $= 0{,}59$). Esistono ragioni strutturali che spiegano perché questo
valore sia un ceiling empirico, non un indicatore di modelli scadenti.

\textbf{Vincolo matematico imposto dall'imbalance.} Con $\approx$9\% di buggy, l'accordo
atteso per caso è $P_e \approx 0{,}836$. Per ottenere $\kappa = 0{,}6$ servirebbe
$P_o = 0{,}934$, cioè il $93{,}4\%$ di predizioni corrette in valore assoluto: ogni
falso positivo e ogni falso negativo incide pesantemente su $P_o$. Raggiungere
simultaneamente alta precision e alto recall con un imbalance così severo è
matematicamente difficile: aumentare il recall genera falsi positivi (abbassa $P_o$),
abbassarlo genera falsi negativi (idem).

\textbf{AUC e Kappa non si contraddicono.} RF ottiene AUC $= 0{,}94$ --- vicino al
perfetto --- ma Kappa $= 0{,}58$. Le due metriche misurano aspetti diversi: AUC valuta
la qualità del \textit{ranking} (threshold-free), mentre Kappa misura la qualità delle
predizioni binarie alla soglia fissa di default ($0{,}5$ in Weka). Con il 9\% di imbalance
la soglia ottimale di classificazione è $\approx 0{,}09$, non $0{,}5$. Il modello sa
discriminare bene (AUC alto), ma la predizione binaria a soglia fissa penalizza Kappa
indipendentemente dalla qualità discriminatoria.

\textbf{Soffitto delle feature statiche.} LOC, churn e commit-count sono correlati ai
bug ma non causalmente collegati: esiste rumore irriducibile che nessun classificatore
può eliminare, ponendo un ceiling empirico sul Kappa raggiungibile con queste feature.

\textbf{Contesto della letteratura.} Hall et al.\ (2012) riportano Kappa tipico
$0{,}2$--$0{,}5$ su 208 studi di defect prediction; Lessmann et al.\ (2008) mostrano che
RF raramente supera $0{,}55$--$0{,}60$ su 10 benchmark. Il valore $0{,}58$--$0{,}59$ di
RF in questo run è nella \textbf{fascia alta della letteratura}.

\subsubsection{Effetto del balancing: trade-off Recall/Kappa}

\begin{table}[h!]
\centering
\small
\begin{tabular}{@{}lccc@{}}
\toprule
\textbf{RF} & \textbf{Precision} & \textbf{Recall} & \textbf{Kappa} \\
\midrule
None          & 0,73 & 0,52 & \textbf{0,58} \\
Undersampling & 0,31 & 0,86 & \textbf{0,37} \\
Oversampling  & 0,61 & 0,66 & \textbf{0,59} \\
\bottomrule
\end{tabular}
\caption{Effetto del balancing su RandomForest: trade-off Precision/Recall/Kappa.}
\end{table}

Undersampling triplica il Recall di RF ($0{,}52 \to 0{,}86$) ma fa crollare il Kappa
($0{,}58 \to 0{,}37$). Lo stesso pattern si osserva su IBk ($0{,}47 \to 0{,}30$).
Undersampling bilancia il training a 50/50, spingendo il classificatore a prevedere
``buggy'' molto più spesso: i falsi negativi calano (Recall alto) ma i falsi positivi
esplodono, facendo crollare la Precision da $0{,}73$ a $0{,}31$ (2 predizioni ``buggy''
su 3 sono errate). Kappa penalizza questa tendenza a sovra-predire la classe positiva
perché $P_o$ cala per via dei numerosi FP.

Oversampling ottiene il miglior bilanciamento complessivo per RF: Precision $0{,}61$,
Recall $0{,}66$, Kappa $0{,}59$. Mantiene un buon Recall senza far esplodere i falsi
positivi, in quanto le istanze minoritarie vengono ricampionate (non si gettano via dati).

\textbf{Interpretazione pratica:} se il costo di un bug sfuggito è altissimo (es.\
sicurezza), Undersampling con alto Recall è preferibile. Se si vuole un bilancio
complessivo, Oversampling su RF è la scelta migliore. Kappa espone una verità che il solo
Recall nasconde: il modello con Undersampling è meno affidabile complessivamente, anche
se cattura più bug.

\subsubsection{AUC di RandomForest stabile al bilanciamento}

\begin{table}[h!]
\centering
\small
\begin{tabular}{@{}lccc@{}}
\toprule
\textbf{Classificatore} & \textbf{AUC (None)} & \textbf{AUC (Under)} & \textbf{AUC (Over)} \\
\midrule
RandomForest & \textbf{0,94} & 0,91 & 0,94 \\
NaiveBayes   & 0,79          & 0,76 & 0,78 \\
IBk          & 0,75          & 0,79 & 0,77 \\
\bottomrule
\end{tabular}
\caption{AUC per classificatore e strategia di balancing.}
\end{table}

Per RF il bilanciamento modifica la soglia effettiva ma non la qualità del ranking: AUC
rimane $0{,}94$ con None e Oversampling, scende leggermente a $0{,}91$ con Undersampling
(per la riduzione dell'82\% del training set). Per IBk l'Undersampling aiuta leggermente
l'AUC ($0{,}75 \to 0{,}79$): il dataset ridotto ($\approx$2.756 istanze) riduce l'effetto
di distorsione della maggioranza nella distanza k-NN.

\subsubsection{NaiveBayes insensibile al bilanciamento}

NB guadagna solo $+3$\,pp di Recall con Undersampling ($0{,}29 \to 0{,}32$), contro i
$+34$\,pp di RF e i $+30$\,pp di IBk. Kappa rimane fisso a $0{,}26$ in tutte e tre le
strategie di balancing.

La causa è strutturale: il bilanciamento modifica la prior $P(\text{buggy})$ nel training,
ma non risolve la violazione dell'assunzione di indipendenza condizionale. Le feature
fortemente correlate (LOC\_Touched, Churn, LOC\_Added) distorcono le likelihood
$P(\text{feature} \mid \text{buggy})$ indipendentemente dalla distribuzione delle classi.
La performance di NB è quindi limitata dall'assumption violation più che dall'imbalance:
SMOTE o FS wrapper difficilmente miglioreranno NB in modo significativo su questo dataset.

\subsubsection{NPofB20: RandomForest eccellente per la prioritizzazione pratica}

\begin{table}[h!]
\centering
\small
\begin{tabular}{@{}lccc@{}}
\toprule
\textbf{Classificatore} & \textbf{NPofB20 (None)} & \textbf{NPofB20 (Under)} &
\textbf{NPofB20 (Over)} \\
\midrule
RandomForest & \textbf{0,88} & 0,80 & 0,87 \\
NaiveBayes   & 0,54          & 0,53 & 0,54 \\
IBk          & 0,56          & 0,57 & 0,61 \\
\bottomrule
\end{tabular}
\caption{NPofB20 per classificatore e strategia di balancing. Baseline casuale: $\approx$0,20.}
\end{table}

RF senza balancing cattura l'\textbf{88\% dei bug nel top 20\% dei file per risk score}:
un team che ispeziona solo il 20\% dei file più a rischio troverebbe quasi tutti i bug.
NB e IBk si fermano a $\approx$0,54--0,61: meglio del casuale (baseline $= 0{,}20$) ma
significativamente inferiori a RF.

La differenza riflette la qualità della calibrazione delle probabilità posterior: RF stima
meglio $P(\text{buggy} \mid \text{features})$ grazie all'averaging su 100 alberi. IBk con
$k=1$ restituisce probabilità binarie (0 o 1), degradando la qualità del ranking; NaiveBayes
tende a polarizzare le probabilità verso gli estremi.

\subsubsection{Analisi delle domande aperte}

Il professore ha posto tre domande analitiche sui risultati:

\begin{enumerate}
\item Quale classificatore è più accurato degli altri?
\item Il miglior classificatore cambia in base a dataset, numero di release e metrica?
\item Qual è il classificatore migliore per una determinata metrica?
\end{enumerate}

Le domande (1) e (3) sono già rispondibili dal CSV di output attuale, aggregando per
\texttt{Classifier} e cercando il massimo per ogni colonna metrica. Dai risultati del run
preliminare: RF domina su AUC e NPofB20; la scelta della strategia di balancing dipende
dalla metrica prioritaria.

La domanda (2) nella parte relativa al \textbf{dataset} richiede l'esecuzione della
pipeline su più progetti (BOOKKEEPER, KAFKA): ogni progetto produce il proprio CSV con la
colonna \texttt{Dataset} valorizzata.

La parte relativa al \textbf{numero di release} richiede invece una modifica
architetturale: la 10$\times$10 CV non ha una nozione di ``quante release sono state usate
per il training'' perché i fold sono casuali. Rispondere richiede la \textbf{Walk-Forward
Validation}, con un training set che cresce release per release e una colonna aggiuntiva
\texttt{NumTrainingReleases} nel CSV di output. L'implementazione richiederebbe un
\texttt{WalkForwardSplitter} che partizioni le istanze rispettando l'ordine temporale
e un \texttt{WalkForwardEvaluatorService} che addestri il classificatore su ogni
training fold e lo valuti sul fold di test (senza cross-validation). Il numero di righe
di output sarebbe: 36 combinazioni $\times$ 10 split $=$ 360 righe. L'implementazione è
ortogonale alla pipeline 10-fold esistente e potrebbe coesistere come modalità alternativa
producendo un secondo CSV.

\subsubsection{Risultati finali}

% TODO: inserire i risultati del run completo (SMOTE abilitato, wrapper FS abilitati,
% dataset aggiuntivi BOOKKEEPER/KAFKA). Sostituire questo placeholder con la tabella
% completa delle 60 combinazioni e le relative analisi comparative.

\medskip
\noindent\textit{I risultati del run completo sono in attesa di completamento.
Questa sezione verrà aggiornata al termine dell'esecuzione.}
